\section{Use Case Diagram}

\subsection*{Attori}
\begin{itemize}
    \item \textbf{Operatore Tecnico:} tecnico del Comune incaricato di costruire scenari, eseguire ottimizzazioni e analizzare i risultati.
    \item \textbf{Decisore Comunale:} soggetto (amministratore, dirigente) che consulta i risultati e le spiegazioni a supporto della decisione finale. Nell'attuale implementazione il ruolo coincide con quello dell'Operatore Tecnico.
    \item \textbf{Assistente AI} (opzionale): componente che produce proposte di configurazione a partire da richieste in linguaggio naturale. È modellato come \emph{attore secondario}.
    \item \textbf{Motore di Ottimizzazione} (CP-SAT): sistema esterno che risolve il problema di ottimizzazione combinatoria. Modellato come \emph{attore secondario}.
    \item \textbf{Servizio Mappa:} sistema esterno che fornisce le tile cartografiche utilizzate per la visualizzazione geografica degli edifici.
\end{itemize}

\subsection*{Diagramma principale}

\begin{lstlisting}[language=,caption={Use Case Diagram principale di TEPO}]
@startuml
left to right direction
skinparam packageStyle rectangle

actor "Operatore Tecnico" as OT
actor "Decisore Comunale" as DC
actor "Assistente AI" as AI <<optional>>
actor "Motore di Ottimizzazione\n(CP-SAT)" as SOLVER <<system>>
actor "Servizio Mappa" as MAP <<system>>

rectangle TEPO {
  usecase "UC-01\nConsultare il catalogo\ndegli edifici" as UC01
  usecase "UC-02\nConsultare il catalogo\ndegli interventi" as UC02
  usecase "UC-03\nConfigurare uno\nscenario" as UC03
  usecase "UC-04\nDefinire i vincoli\ndello scenario" as UC04
  usecase "UC-05\nEseguire\nl'ottimizzazione" as UC05
  usecase "UC-06\nConsultare risultati\ne spiegazioni" as UC06
  usecase "UC-07\nSalvare e confrontare\nscenari A/B" as UC07
  usecase "UC-08\nAnalisi di\nsensibilita sul budget" as UC08
  usecase "UC-09\nOttenere una proposta\nassistita dall'AI" as UC09
  usecase "UC-10\nApprovare o rifiutare\nla proposta AI" as UC10
  usecase "UC-11\nImportare dati\nda CSV" as UC11
  usecase "UC-12\nConsultare lo storico\ndelle ottimizzazioni" as UC12
  usecase "UC-13\nVisualizzare gli edifici\nsu mappa" as UC13
}

OT --> UC01
OT --> UC02
OT --> UC03
OT --> UC04
OT --> UC05
OT --> UC06
OT --> UC07
OT --> UC08
OT --> UC09
OT --> UC10
OT --> UC11
OT --> UC12
OT --> UC13
DC --> UC06
DC --> UC07

UC05 ..> SOLVER : <<include>>
UC13 ..> MAP : <<include>>
UC09 ..> AI : <<include>>
UC10 .> UC09 : <<extend>>
UC07 ..> UC06 : <<include>>
UC08 ..> UC05 : <<include>>

@enduml
\end{lstlisting}

% IMMAGINE: img/UseCaseDiagram.png -- Diagramma dei casi d'uso principale di TEPO

\subsection{Descrizione dei casi d'uso principali}

\paragraph{\large UC-03: Configurare uno scenario di ottimizzazione}

\paragraph{RF collegati:} RF\ref{rf:config_scenario}, RF\ref{rf:candidati}, RF\ref{rf:vincoli_scenario}.

\paragraph{Attore primario:} Operatore Tecnico.

\paragraph{Obiettivo:} definire in modo completo i parametri (budget, obiettivo, candidati, vincoli) di uno scenario pronto per l'esecuzione dell'ottimizzazione.

\paragraph{Precondizioni:} il catalogo degli edifici e degli interventi è popolato; l'utente ha accesso all'interfaccia di configurazione dello scenario.

\paragraph{Postcondizioni:} lo scenario corrente contiene un budget valido, un obiettivo selezionato, un insieme non vuoto di candidati e un insieme di vincoli ben formati.

\paragraph{Flusso principale:}
\begin{enumerate}
    \item L'utente apre la pagina di configurazione dello scenario.
    \item Il sistema mostra i parametri correnti dello scenario (budget, obiettivo, candidati, vincoli).
    \item L'utente inserisce o modifica il budget complessivo.
    \item L'utente seleziona l'obiettivo tra risparmio energetico, risparmio economico e riduzione di CO\textsubscript{2}.
    \item L'utente seleziona le opzioni di intervento candidate a partire dal catalogo. [Estensione 1]
    \item L'utente aggiunge eventuali vincoli di incompatibilità, prerequisito, numero massimo, inclusione o esclusione. [Estensione 2]
    \item Il sistema verifica la coerenza logica dei parametri. [Eccezione 1]
    \item Il sistema conferma all'utente che lo scenario è pronto per l'ottimizzazione.
\end{enumerate}

\paragraph{Estensioni:}
\begin{enumerate}
    \item L'utente può utilizzare le funzioni di selezione rapida (seleziona tutti / deseleziona tutti) oppure filtrare il catalogo per edificio e tipologia prima di selezionare i candidati.
    \item L'utente può rimuovere un vincolo precedentemente aggiunto.
\end{enumerate}

\paragraph{Eccezioni:}
\begin{enumerate}
    \item Se il budget è nullo o negativo, se non è stato selezionato alcun obiettivo, oppure se la stessa opzione risulta contemporaneamente inclusa ed esclusa, il sistema blocca il salvataggio e mostra un messaggio di errore specifico.
\end{enumerate}

\paragraph{\large UC-05: Eseguire l'ottimizzazione e consultare le spiegazioni}

\paragraph{RF collegati:} RF\ref{rf:ottimizzazione}, RF\ref{rf:risultati}, RF\ref{rf:spiegazioni}.

\paragraph{Attore primario:} Operatore Tecnico.

\paragraph{Attori secondari:} Motore di Ottimizzazione.

\paragraph{Obiettivo:} ottenere il portafoglio ottimo di interventi per lo scenario corrente e comprenderne le motivazioni.

\paragraph{Precondizioni:} lo scenario corrente è completo e coerente (UC-03 concluso con successo).

\paragraph{Postcondizioni:} il sistema mostra il portafoglio ottimo, gli indicatori aggregati e le spiegazioni per ciascuna opzione candidata. L'esecuzione è registrata nello storico delle ottimizzazioni.

\paragraph{Flusso principale:}
\begin{enumerate}
    \item L'utente richiede l'esecuzione dell'ottimizzazione.
    \item Il sistema valida nuovamente lo scenario.
    \item Il sistema invoca il motore di ottimizzazione. [Estensione 1]
    \item Il motore di ottimizzazione restituisce una soluzione con lo stato del calcolo.
    \item Il sistema calcola le spiegazioni deterministiche per ciascuna opzione candidata.
    \item Il sistema registra l'esecuzione nello storico ai fini di audit.
    \item Il sistema presenta all'utente il portafoglio ottimo, gli indicatori aggregati e le spiegazioni.
\end{enumerate}

\paragraph{Estensioni:}
\begin{enumerate}
    \item In caso di più soluzioni ottime equivalenti, il sistema applica criteri di tie-break deterministici (minimizzazione del costo, minimizzazione del numero di interventi, preferenza per identificativi crescenti) e ne dà evidenza all'utente.
\end{enumerate}

\paragraph{Eccezioni:}
\begin{enumerate}
    \item Se il problema risulta inammissibile, il sistema mostra lo stato ``inammissibile'' e propone all'utente una diagnostica sulle possibili cause (es. vincolo di inclusione su un'opzione non candidata).
    \item Se il calcolo supera il tempo massimo, il sistema mostra lo stato ``indeterminato per timeout'' e non presenta alcun portafoglio come definitivo.
\end{enumerate}

\paragraph{\large UC-07: Salvare e confrontare due scenari}

\paragraph{RF collegati:} RF\ref{rf:scenari_salvati}, RF\ref{rf:confronto}.

\paragraph{Attore primario:} Operatore Tecnico.

\paragraph{Attore secondario:} Decisore Comunale.

\paragraph{Obiettivo:} confrontare in modo strutturato due configurazioni alternative (scenario A e scenario B) per orientare la decisione.

\paragraph{Precondizioni:} almeno due scenari sono stati configurati ed eventualmente ottimizzati.

\paragraph{Postcondizioni:} il sistema mostra in modo comparato i parametri e gli indicatori aggregati dei due scenari, evidenziando gli scostamenti.

\paragraph{Flusso principale:}
\begin{enumerate}
    \item L'utente salva lo scenario corrente nello slot A.
    \item L'utente modifica lo scenario corrente e lo salva nello slot B. [Estensione 1]
    \item L'utente accede alla pagina di confronto.
    \item Il sistema recupera dagli slot A e B i parametri e gli indicatori aggregati delle rispettive esecuzioni.
    \item Il sistema mostra i due scenari in forma tabellare, evidenziando gli scostamenti.
\end{enumerate}

\paragraph{Estensioni:}
\begin{enumerate}
    \item L'utente può sovrascrivere lo slot A o lo slot B con lo scenario corrente in qualunque momento, dopo conferma esplicita.
\end{enumerate}

\paragraph{\large UC-09: Ottenere una proposta assistita dall'AI}

\paragraph{RF collegati:} RF\ref{rf:ai_proposta}, RF\ref{rf:ai_approvazione}.

\paragraph{Attore primario:} Operatore Tecnico.

\paragraph{Attore secondario:} Assistente AI (opzionale).

\paragraph{Obiettivo:} ottenere una proposta preliminare di configurazione dello scenario a partire da una richiesta in linguaggio naturale.

\paragraph{Precondizioni:} l'utente si trova nella pagina di configurazione dello scenario.

\paragraph{Postcondizioni:} il sistema mostra una proposta di configurazione corredata da origine, eventuali segnalazioni e richiesta esplicita di approvazione umana.

\paragraph{Flusso principale:}
\begin{enumerate}
    \item L'utente apre il pannello di assistenza AI.
    \item L'utente digita la propria richiesta in linguaggio naturale.
    \item Il sistema inoltra la richiesta al componente di assistenza (motore a regole oppure, se configurato, modello linguistico esterno). [Estensione 1]
    \item Il componente di assistenza restituisce una proposta di configurazione con origine, eventuali informazioni mancanti ed eventuali avvisi.
    \item Il sistema mostra la proposta all'utente senza applicarla.
    \item L'utente accetta, modifica oppure rifiuta la proposta. [Eccezione 1]
    \item In caso di accettazione, il sistema applica i parametri proposti allo scenario corrente e ne dà evidenza all'utente.
\end{enumerate}

\paragraph{Estensioni:}
\begin{enumerate}
    \item Se il modello linguistico esterno non è raggiungibile o restituisce una risposta non conforme, il sistema ricade sul motore a regole e segnala l'evento all'utente.
\end{enumerate}

\paragraph{Eccezioni:}
\begin{enumerate}
    \item Se l'utente rifiuta la proposta, il sistema non modifica lo scenario corrente.
\end{enumerate}